\documentclass[11pt]{article}

\usepackage[margin=25mm]{geometry}

\usepackage{amsmath,amssymb,fontspec,unicode-math}

\IfFontExistsTF{DejaVu Serif}{\setmainfont{DejaVu Serif}}{\setmainfont{Latin Modern Roman}}

\IfFontExistsTF{Latin Modern Math}{\setmathfont{Latin Modern Math}}{\setmathfont{DejaVu Math TeX Gyre}}

\usepackage{xurl,hyperref}

\setlength{\parindent}{0pt}\setlength{\parskip}{6pt}

\begin{document}

\title{Breaking the Discretization Barrier of Continuous Physics Simulation Learning}

\author{ELVIA preliminary AI analysis}\date{21 October 2026}\maketitle

Preliminary assessment: Minor.\par

Independent scientific review remains pending. Author adoption has not been requested.

The post-correction error bound is valid under its stated recurrence, but it does not give the same bound between corrections or force actual error to converge. Additive resets and cell-dependent graph neighborhoods also do not automatically preserve temporal or spatial continuity.

\section*{Review scope}

Claim-focused reading of the original paper and relevant author-released passages; direct mathematical or static-source reasoning. No author program, earlier verifier, model or benchmark was executed during this check.

\section*{Problem 1: The geometric error bound needs a post-correction domain}

Source: Theorem 3.1, Eqs.(15)–(16), PDF p.6; Appendix F, Eqs.(24)–(27), p.24

Assessment: supported; impact: minor

\subsection*{Source claim (paraphrase)}

The main display uses \(K=\lfloor(t-t_0)/\Delta\rfloor\) for time \(t\), while the proof establishes a bound on error just after correction \(K\).

\subsection*{Analysis}

Write \(e_k^+\) for corrected error. The assumptions imply \[e_{k+1}^+\le\alpha e_k^++B,\quad \alpha=\kappa e^{L_F\Delta}<1,\quad B=\kappa E_{\rm ODE}+\delta_C.\] Hence \(e_K^+\le\alpha^K e_0^++B(1-\alpha^K)/(1-\alpha)\). This is the valid recurrence proved in Appendix F.

An arbitrary-time extension needs an unfinished-flow factor. For ideal trajectory \(0\), exact learned ODE \(\dot y=y\), \(\Delta=1\), initial corrected value \(1\) and correction \(y^+=e^{-2}y^-\), the assumptions hold with \(L_F=1,\kappa=e^{-2},E_{\rm ODE}=\delta_C=0\). Post-correction errors are \(e^{-k}\), but at \(t=1/2\) the error is \(e^{1/2}>1\), exceeding Eq.(16) read literally with \(K=0\).

\subsection*{Finding and scope}

The proof’s explicit post-correction interpretation resolves the principal domain issue. The same floor-indexed bound cannot be used unchanged at arbitrary times between corrections.

\subsection*{Proposed correction}

Label the bound as applying to \(e_K^+\). If a partial-interval residual \(E_{\rm ODE}(\tau)\) is established, use \(e(t_K+\tau)\le e^{L_F\tau}e_K^++E_{\rm ODE}(\tau)\).

A proposed correction is retained where stated; independent verification and author adoption remain unestablished.

\section*{Problem 2: A convergent upper envelope does not force error convergence}

Source: Theorem 3.1 after Eq.(16), PDF p.6; final paragraph of Appendix F, p.24

Assessment: supported; impact: minor

\subsection*{Source claim (paraphrase)}

The paper concludes that actual corrected error converges to the nonzero floor \(C_2=B/(1-\alpha)\).

\subsection*{Analysis}

The inequality establishes only \(\limsup_K e_K^+\le C_2\). For a smooth scalar counterexample take ideal state \(0\), zero learned ODE, \(\Delta=1\), initial state \(0\), and corrected-state map \(G(y)=\cos^2(\pi y/2)\). Its additive corrector is \(G(y)-y\). Since \(|G(y)|\le1\le |y|/2+1\), assumption (b) holds with \(\kappa=1/2,\delta_C=1\); take \(E_{\rm ODE}=0\), so \(\alpha=1/2,C_2=2\). The trajectory alternates \(0,1,0,1,\ldots\), never converging, while satisfying every displayed upper bound. The assumption controls error relative to the ideal state; it does not require pairwise contraction of \(G\).

\subsection*{Finding and scope}

Bounded long-horizon error follows conditionally, but existence of an error limit or equality to the envelope floor does not.

\subsection*{Proposed correction}

Replace the limit claim by \(\limsup_K e_K^+\le C_2\). If \(B=0\), the existing inequality does imply convergence to zero; otherwise extra dynamical assumptions are needed.

A proposed correction is retained where stated; independent verification and author adoption remain unestablished.

\section*{Problem 3: Nonzero additive correction creates a temporal jump}

Source: Section 3.2, Eqs.(11)–(12) and final inference paragraph, PDF p.5

Assessment: supported; impact: minor

\subsection*{Source claim (paraphrase)}

The hybrid inference strategy is said to maintain continuity and smoothness although it adds a discrete correction after each ODE interval.

\subsection*{Analysis}

Equation (12) uses the previous corrected state in a lagged residual. Let the ODE be zero and initial state \(1\); choose \(\lambda=1/2\) and an admissible constant residual \(r_\psi(1)=-1/2\). The state immediately before reset is \(1\), while immediately after it is \(3/4\). The jump is \(-1/4\). In general continuity requires the added residual to vanish at the reset. The theorem explicitly permits pre/post states, so this finding does not contradict its hybrid formulation.

\subsection*{Finding and scope}

Arbitrary-time queries and piecewise smooth ODE segments remain possible; globally continuous trajectories are not guaranteed by the printed reset construction.

\subsection*{Proposed correction}

Describe a piecewise smooth hybrid trajectory and state which side of a reset defines a query. A continuously forced replacement would be a different method requiring its own analysis.

A proposed correction is retained where stated; independent verification and author adoption remain unestablished.

\section*{Problem 4: A changing four-vertex neighborhood can break spatial continuity}

Source: Section 3.1, Eq.(4), PDF p.4; Section 3.3, pp.5–6; Algorithm 1, p.22; Appendix C.1, p.23

Assessment: supported; impact: minor

\subsection*{Source claim (paraphrase)}

Relative-position message passing over the four vertices of the query’s cell is used to support spatial continuity.

\subsection*{Analysis}

Consider adjacent cells \([0,1]\times[0,1]\) and \([1,2]\times[0,1]\). Set all vertex features at \(x=0,1\) to zero and both at \(x=2\) to one; these values can sample the smooth field \(x(x-1)/2\). Choose zero query Gabor initialization, zero positional term and bias, and first message weight \(1\). As a query crosses \(x=1\), Eq.(4) sums features of four vertices: its one-sided outputs are \(0\) and \(2\). A second residual layer with zero weight/bias is the identity, and a continuous ReLU readout preserves this nonnegative jump. Within each fixed cell the map can be continuous; neighbor-set switching is the obstruction.

\subsection*{Finding and scope}

The printed decoder does not guarantee continuity across cell boundaries. Extra interpolation in an unavailable implementation could change this conclusion, and no trained spatial discontinuity is asserted.

\subsection*{Proposed correction}

A shared-vertex bilinear interpolation construction would enforce continuity across cell boundaries, but not necessarily differentiability. This proposed decoder change has not been validated as preserving trained performance.

A proposed correction is retained where stated; independent verification and author adoption remain unestablished.

\section*{Limits}

Preliminary AI assessment; independent human scientific adjudication remains pending.\par

Analytic witnesses and static release checks do not remeasure the published benchmark results.\par

The published author repository could not be resolved in the current metadata check; conclusions about architecture are bounded to the paper’s printed specification.\par

\section*{Source identity}

NeurIPS 2025 official proceedings PDF (retained local copy)

\url{https://proceedings.neurips.cc/paper_files/paper/2025/file/74b7956113fdf0ec87288f351a1d8a34-Paper-Conference.pdf}

Reviewed PDF SHA-256: \nolinkurl{85fcfe46da8155f23ee8bebb6e927b9d6ccbed2b82b2611d530d7d7bf20a87f5}

The SHA-256 identifies the exact retained PDF inspected in this rapid source/argument check. Original supplied report bytes are retained locally. No unambiguous original PDF retrieval timestamp was recovered; file modification time was not substituted.

\end{document}
